% !TeX root = ../main.tex
\section{FRI: Fast Reed--Solomon IOPs of Proximity}
\label{sec:fri}

The polynomial IOP chapters above repeatedly used the same verifier strategy: commit to a
polynomial, evaluate it at a random point, and use Schwartz--Zippel to detect a false identity.
This strategy presupposes that the committed oracle really is the evaluation of a
\emph{low-degree} polynomial. KZG enforces this through a structured reference string: the
commitment can only be formed from the published powers up to the supported degree. A
hash-based system has no such algebraic restriction. A Merkle root can bind the prover to an
arbitrary table, including a table that is not close to any low-degree polynomial.

FRI fills precisely this gap~\cite{BenSassonBentovHoreshRiabzev2018FRI}. Its full name is
the Fast Reed--Solomon Interactive Oracle Proof of Proximity. It lets a verifier with oracle
access to a large evaluation table test, using only a small number of queries, whether that
table is close to a Reed--Solomon codeword. Its central operation resembles one layer of an inverse FFT:
split a polynomial into its even and odd coefficient polynomials, combine those two pieces
using a random verifier challenge, and recurse on a domain of half the size.

The main conceptual point is:
\[
\boxed{
\begin{array}{c}
\text{Merkle commitment binds an arbitrary evaluation table,}\\[0.2em]
\text{FRI proves that the table is close to a low-degree polynomial.}
\end{array}}
\]
Neither component replaces the other. Together with Fiat--Shamir, they form the low-degree
commitment layer used by many transparent proof systems.

\subsection{The Reed--Solomon proximity problem}

Let $D\subseteq\F$ be a set of $N$ distinct points and let $1\le d\le N$. Define the
degree-$<d$ Reed--Solomon code on $D$ by
\[
    \RS[\F,D,d]
    =\set{\bigl(p(x)\bigr)_{x\in D}:p\in\F[X],\ \deg p<d}
    \subseteq \F^D.
\]
The dimension is $d$ and the rate is
\[
    \rho=\frac{d}{N}.
\]
The reciprocal
\[
    B=\rho^{-1}=\frac{N}{d}
\]
is the \emph{FRI blowup factor}. Starting from the $d$ coefficients of a degree-$<d$
polynomial, the prover performs a low-degree extension to $N=Bd$ evaluations and commits
only to those evaluations, rather than to all $|\F|$ field points. Thus $B$ exposes the main
engineering tradeoff:
\begin{itemize}[leftmargin=2em]
    \item increasing $B$ increases the prover's evaluation and hashing work, and the initial
    Merkle tree, essentially linearly;
    \item increasing $B$ lowers the code rate and increases the distance, so fewer random
    query paths are usually needed for the same target soundness.
\end{itemize}
At the lecture-level heuristic where a far word has disagreement rate near $1-\rho$, one
query misses the disagreement with probability about $\rho=1/B$. Hence $t$ independent
paths contribute at most roughly
\[
    t\log_2 B
\]
bits of query soundness, suggesting $t\approx\lceil\lambda/\log_2B\rceil$ for a
$\lambda$-bit target. This is a parameter-selection heuristic, not a complete FRI security
bound: finite-field terms, list-decoding radii, the exact FRI variant, and Fiat--Shamir
grinding must also be included.

If $p$ and $q$ are distinct polynomials of degree less than $d$, then $p-q$ has at most
$d-1$ roots. Their evaluation vectors therefore disagree in at least $N-d+1$ positions.
Consequently, the relative minimum distance of the code is
\[
    \delta_{\min}=\frac{N-d+1}{N}>1-\rho.
\]

For two functions $f,g:D\to\F$, their relative Hamming distance is
\[
    \Delta_D(f,g)
    =\frac{1}{N}\abs{\set{x\in D:f(x)\neq g(x)}}.
\]
The distance from $f$ to the code is
\[
    \Delta_D\bigl(f,\RS[\F,D,d]\bigr)
    =\min_{\deg p<d}\Delta_D(f,p|_D).
\]

\begin{definition}[Reed--Solomon IOP of proximity]
An interactive oracle proof of proximity for $\RS[\F,D,d]$ gives the verifier oracle access
to a claimed word $f^{(0)}:D\to\F$ and to auxiliary prover oracles. It has:
\begin{itemize}[leftmargin=2em]
    \item \textbf{perfect completeness} if every $f^{(0)}\in\RS[\F,D,d]$ is accepted by
    the honest verifier;
    \item \textbf{proximity soundness} $s(\delta)$ if every word at relative distance
    $\delta$ from the code is rejected with probability at least $s(\delta)$, regardless of
    the auxiliary oracles supplied by a cheating prover.
\end{itemize}
\end{definition}

The word \emph{proximity} matters. A verifier making $q\ll N$ queries cannot reject every
word that differs from a codeword in only one position: it may simply miss that position.
FRI therefore gives a rejection probability that grows with the distance from the code. It is
not a deterministic proof of exact code membership.

\subsection{Smooth multiplicative domains}

We use the multiplicative presentation most compatible with the preceding chapters. Let
$\F$ have odd characteristic and suppose that $\F^\times$ contains a subgroup of order
\[
    N=2^m.
\]
Let $\omega$ generate that subgroup and let $\gamma\in\F^\times$. Define
\[
    D_0=\gamma\angles{\omega}.
\]
Successively square the domains:
\[
    D_i
    =\set{x^{2^i}:x\in D_0}
    =\gamma^{2^i}\angles{\omega^{2^i}},
    \qquad
    |D_i|=\frac{N}{2^i}.
\]
For every $i<m$, the squaring map
\[
    q_i:D_i\to D_{i+1},
    \qquad
    q_i(x)=x^2,
\]
is two-to-one. The two preimages of $y=x^2$ are $x$ and $-x$.

In a prime field $\F_p$, a subgroup of order $N$ exists exactly when $N\mid(p-1)$. This
divisibility requirement explains the use of FFT-friendly fields. For example,
\[
    p=2^{64}-2^{32}+1,
    \qquad
    p-1=2^{32}(2^{32}-1),
\]
supports power-of-two subgroups through order $2^{32}$. Write
$\gamma_i=\gamma^{2^i}$, $\omega_i=\omega^{2^i}$, and $N_i=|D_i|$. Since $N_i$ is even,
\[
    -\gamma_i\omega_i^j=\gamma_i\omega_i^{j+N_i/2}.
\]
Thus the two entries paired by a binary fold have indices separated by $N_i/2$, and both map
to the same point of $D_{i+1}$. Implementations commonly place the two values in the same
Merkle leaf, or at least adjacent in the leaf ordering, so that opening a fold exposes and
authenticates the pair together.

For clarity, assume first that the degree bound is a power of two,
\[
    d_0=2^r,
    \qquad r\le m,
\]
and set
\[
    d_i=\frac{d_0}{2^i}.
\]
Both the domain size and the degree bound halve in every round, so the code rate remains
constant:
\[
    \frac{d_i}{|D_i|}=\frac{d_0}{N}=\rho.
\]
After $r$ rounds the alleged polynomial has degree less than $1$, so it must be constant.

\begin{remark}[Relation to the original FRI presentation]
The formal main theorem of the original FRI paper is stated for additive cosets in binary
extension fields. In characteristic two, $x$ and $-x$ coincide, so the map $x\mapsto x^2$
is not the required two-to-one map; the construction instead uses affine-subspace
polynomials, which are linearized polynomials. The same paper explains the smooth
multiplicative-group version used here. Modern expositions over prime fields commonly use
this roots-of-unity formulation because it exposes the FFT-style even/odd decomposition
directly~\cite{BenSassonBentovHoreshRiabzev2018FRI}.
\end{remark}

\subsection{One folding round}

Let
\[
    p(X)=a_0+a_1X+\cdots+a_{d-1}X^{d-1}
\]
have degree less than $d$. Separate its even and odd coefficients:
\[
    p_{\mathrm{even}}(Y)=\sum_{j\ge 0}a_{2j}Y^j,
    \qquad
    p_{\mathrm{odd}}(Y)=\sum_{j\ge 0}a_{2j+1}Y^j.
\]
Then
\[
\boxed{
    p(X)=p_{\mathrm{even}}(X^2)+X\,p_{\mathrm{odd}}(X^2).
}
\]
Both component polynomials have degree less than $\lceil d/2\rceil$. Given a verifier
challenge $\alpha\in\F$, define the folded polynomial
\[
    \operatorname{Fold}_{\alpha}(p)(Y)
    =p_{\mathrm{even}}(Y)+\alpha p_{\mathrm{odd}}(Y).
\]
Therefore
\[
    \deg \operatorname{Fold}_{\alpha}(p)<\left\lceil\frac d2\right\rceil.
\]
This is the source of FRI's degree reduction.

The verifier does not know the coefficients of $p$. It only has oracle access to values
$f(x)=p(x)$ on $D_i$. Nevertheless, the folded value at $y=x^2$ can be recovered from the
pair $x,-x$:
\[
    p(x)=p_{\mathrm{even}}(y)+x p_{\mathrm{odd}}(y),
\]
\[
    p(-x)=p_{\mathrm{even}}(y)-x p_{\mathrm{odd}}(y).
\]
Since the characteristic is odd,
\[
    p_{\mathrm{even}}(y)=\frac{p(x)+p(-x)}{2},
\]
\[
    p_{\mathrm{odd}}(y)=\frac{p(x)-p(-x)}{2x}.
\]
Hence the round-consistency equation is
\begin{tcolorbox}
\[
\boxed{
f^{(i+1)}(x^2)
=
\frac{f^{(i)}(x)+f^{(i)}(-x)}{2}
+\alpha_i\frac{f^{(i)}(x)-f^{(i)}(-x)}{2x}.
}
\tag{FRI-fold}
\]
\end{tcolorbox}

Equivalently, collecting the coefficients of the two opened values gives the Lagrange-weight
form emphasized in many implementations:
\[
\boxed{
f^{(i+1)}(x^2)
=
\frac{\alpha_i+x}{2x}f^{(i)}(x)
+\frac{\alpha_i-x}{-2x}f^{(i)}(-x).
}
\]
The two weights sum to one. As a useful sanity check, setting $\alpha_i=x$ returns
$f^{(i)}(x)$, while setting $\alpha_i=-x$ returns $f^{(i)}(-x)$.

There is a useful interpolation interpretation. For fixed $y=x^2$, let $\ell_y(Z)$ be the
unique affine polynomial satisfying
\[
    \ell_y(x)=f^{(i)}(x),
    \qquad
    \ell_y(-x)=f^{(i)}(-x).
\]
Then
\[
    \ell_y(Z)
    =\frac{f^{(i)}(x)+f^{(i)}(-x)}{2}
     +Z\frac{f^{(i)}(x)-f^{(i)}(-x)}{2x},
\]
and the folding check is simply
\[
    f^{(i+1)}(y)\stackrel{?}{=}\ell_y(\alpha_i).
\]
Thus each fiber $\{x,-x\}$ defines a line, and the verifier asks the prover to supply the
values of all those lines at one unpredictable horizontal coordinate $\alpha_i$.

\begin{intuitionbox}[title=\textbf{Why the challenge comes after the current oracle}]
If the prover could see $\alpha_i$ before binding itself to $f^{(i)}$, it could choose the two
values in every pair so that their interpolating line passes through any desired
$f^{(i+1)}(x^2)$. The order
\[
    \text{commit to }f^{(i)}
    \longrightarrow \alpha_i
    \longrightarrow \text{commit to }f^{(i+1)}
\]
prevents this backwards construction. The random linear combination plays the same role as
the verifier challenge in the sigma protocols studied earlier.
\end{intuitionbox}

\subsection{The pedagogical binary FRI protocol}

We now state the complete arity-two protocol. The word $f^{(0)}:D_0\to\F$ is claimed to be
close to $\RS[\F,D_0,d_0]$, where $d_0=2^r$. In the IOP model, ``send an oracle'' means
that the verifier may later query selected entries without reading the whole message.

\subsubsection*{Commit phase}

The verifier begins with oracle access to $f^{(0)}$. For $i=0,\ldots,r-1$:
\begin{enumerate}[leftmargin=2em]
    \item the verifier samples a uniform challenge $\alpha_i\leftarrow\F$;
    \item the prover sends a new oracle
    \[
        f^{(i+1)}:D_{i+1}\to\F,
    \]
    allegedly satisfying the folding equation (FRI-fold) for every $x\in D_i$.
\end{enumerate}
The final oracle $f^{(r)}$ is alleged to be constant. The prover sends that constant
$c\in\F$. More generally, an implementation may stop earlier and send the coefficients of a
small terminal polynomial, which the verifier evaluates directly.

The final domain has not generally collapsed to one point. Since $r=\log_2d_0$,
\[
    |D_r|=\frac{N}{d_0}=B.
\]
It is the \emph{degree bound} that has collapsed to zero. Sending one field element $c$ is
therefore an assertion that all $B$ remaining evaluations equal the same constant. If the
protocol stops after $s<r$ folds, the terminal polynomial instead has degree
$<d_0/2^s$; the prover sends its coefficients, and the verifier evaluates that polynomial at
the terminal point reached by each query path.

The honest prover does not need to interpolate a polynomial in every round. For each pair
$\{x,-x\}$, it computes one output value using (FRI-fold). The total number of field
operations over all rounds is therefore a geometric sum:
\[
    O(N)+O(N/2)+O(N/4)+\cdots=O(N).
\]

\subsubsection*{Query phase}

After all oracles have been fixed, the verifier samples $t$ independent query seeds
\[
    x_0^{(1)},\ldots,x_0^{(t)}\leftarrow D_0.
\]
For each path $j$ set recursively
\[
    x_{i+1}^{(j)}=\bigl(x_i^{(j)}\bigr)^2.
\]
At round $i$, the verifier queries
\[
    f^{(i)}\bigl(x_i^{(j)}\bigr),
    \qquad
    f^{(i)}\bigl(-x_i^{(j)}\bigr),
    \qquad
    f^{(i+1)}\bigl(x_{i+1}^{(j)}\bigr),
\]
and checks (FRI-fold). The value in the third query is reused as one of the two values in the
next round. In the last round it is compared with the terminal constant $c$.

The verifier accepts only if every folding equation on every path is valid and the terminal
degree check succeeds.

\begin{tcolorbox}[colback=blue!4!white,colframe=blue!45!black]
\textbf{Commit-before-query discipline.} The query locations must be sampled only after all
oracles have been fixed. Otherwise the prover could make the checked paths consistent and
place arbitrary values everywhere else. In the non-interactive version, this ordering is
enforced by deriving every challenge and query seed from the transcript prefix that precedes
it.
\end{tcolorbox}

\subsection{Perfect completeness}

\begin{lemma}[Degree reduction under folding]
If $\deg p_i<d_i$, then for every $\alpha_i\in\F$,
\[
    p_{i+1}=\operatorname{Fold}_{\alpha_i}(p_i)
\]
satisfies
\[
    \deg p_{i+1}<\left\lceil\frac{d_i}{2}\right\rceil.
\]
If $d_i$ is even, this is $\deg p_{i+1}<d_i/2$.
\end{lemma}

\begin{proof}
Writing
\[
    p_i(X)=p_{i,\mathrm{even}}(X^2)+X p_{i,\mathrm{odd}}(X^2),
\]
the largest exponent appearing in either component is obtained by dividing an exponent of
$p_i$ by two and rounding down. Thus both component degrees are less than
$\lceil d_i/2\rceil$. Their linear combination has no larger degree.
\end{proof}

\begin{theorem}[Perfect completeness of binary folding]
Suppose $f^{(0)}=p_0|_{D_0}$ for some polynomial $p_0$ of degree less than $d_0=2^r$.
If the prover sets
\[
    p_{i+1}=\operatorname{Fold}_{\alpha_i}(p_i),
    \qquad
    f^{(i+1)}=p_{i+1}|_{D_{i+1}},
\]
then every verifier check passes, and $p_r$ is constant.
\end{theorem}

\begin{proof}
The folding equation follows from the even/odd decomposition, so every local query check is
an identity. By the lemma and induction,
\[
    \deg p_i<\frac{d_0}{2^i}.
\]
At $i=r$ this gives $\deg p_r<1$, hence $p_r$ is constant and equals the value sent by the
honest prover.
\end{proof}

\subsection{Why local consistency implies proximity soundness}

Completeness is elementary; soundness is the deep part of FRI. A local folding check only
shows consistency among three queried values. A cheating prover may choose each oracle to be
an arbitrary function, not the evaluation of any polynomial. The analysis must show that a
word far from the initial Reed--Solomon code cannot be threaded through all folding rounds
while making most local checks pass.

For fixed oracles and a fixed challenge, define the round inconsistency rate
\[
\begin{aligned}
\varepsilon_i
=\Pr_{x\leftarrow D_i}\Biggl[
f^{(i+1)}(x^2)\neq{}&
\frac{f^{(i)}(x)+f^{(i)}(-x)}{2}\\
&+\alpha_i\frac{f^{(i)}(x)-f^{(i)}(-x)}{2x}
\Biggr].
\end{aligned}
\]
If $\varepsilon_i$ is large, a random query catches the inconsistency directly. If it is
small, then $f^{(i+1)}$ agrees on most points with the random fold of $f^{(i)}$. The crucial
FRI proximity-gap theorem says, roughly, that random folding does not usually turn a word
that is far from a Reed--Solomon code into one that is much closer to the smaller code.

The reason randomness is necessary can be seen from list decoding. A received word may be
close to several low-degree candidates. Each candidate pair induces algebraic conditions on
$\alpha_i$ under which their folded images collide or become unusually close. Because the
list of relevant candidates is controlled and $\alpha_i$ is sampled after the current oracle
is fixed, only a small set of challenges is bad. Outside that set, either distance survives in
the next round or the prover creates noticeable local inconsistency.

The simplified security calculation presented in many lectures separates two failure modes as
\[
    \Pr[\text{accept}]
    \lesssim
    \varepsilon_{\mathrm{field}}+(1-\delta)^t.
\]
Here $\varepsilon_{\mathrm{field}}$ accounts for exceptional folding challenges---in a basic
Schwartz--Zippel-style discussion it is on the order of a degree bound divided by $|\F|$---and
$(1-\delta)^t$ is the probability that all $t$ query paths miss a $\delta$ fraction of bad
locations. This formula is valuable intuition, but it is not a substitute for the theorem of
the particular FRI variant and proximity regime being instantiated.

There is also a matching interpolation attack explaining why the rate cannot be ignored.
Assume for simplicity that $d$ is even. A cheating prover can choose a union $T$ of $d/2$
pairs $\{x,-x\}\subset D_0$, so $|T|=d=\rho N$, and interpolate the unique degree-$<d$
polynomial $s$ agreeing with the initial word on $T$. It then builds every later oracle by
honestly folding $s$, not the full initial word. A query path whose first pair lies in $T$
sees a perfectly consistent transcript; a uniformly sampled pair lies in $T$ with probability
$\rho$. Consequently, all $t$ paths pass with probability
\[
    \rho^t=2^{-t\log_2(1/\rho)}.
\]
This attack works even when the initial word is otherwise very far from the code. It shows that
no analysis can credit a binary-FRI path with more than about $\log_2(1/\rho)$ bits merely from
sampling. It also explains why a larger blowup factor improves query efficiency.

The original FRI theorem makes this quantitative for Reed--Solomon codes of rate
$\rho=2^{-R}$, $R\ge2$. For a word at distance $\delta$ from the code, its optimized IOPP
rejects with probability at least
\[
    \min\set{\delta(1-o(1)),\delta_0},
\]
where $\delta_0>0$ depends mainly on the code rate. For block length $N$, the construction
has linear prover arithmetic, logarithmic verifier arithmetic, and logarithmic query
complexity~\cite{BenSassonBentovHoreshRiabzev2018FRI}. The precise theorem uses the
additive-domain, higher-arity protocol rather than the pedagogical binary version above, but
the same proximity-preserving mechanism is at work.

The many-round structure also calls for a stronger, round-by-round view of soundness. Later work
proved round-by-round soundness and knowledge soundness for FRI and Batched FRI, and used those
properties to justify a non-interactive random-oracle compilation
\cite{BlockGarretaKatzThalerTiwariZajac2023FRIFS}. A representative bound in the regime
$\delta<1-\sqrt{\rho}$ has the schematic form
\[
    \varepsilon_{\mathrm{rbr}}
    =\max\set{
      O\!\left(\frac{N^2}{\rho^{3/2}|\F|}\right),
      (1-\delta)^t
    },
\]
with exact notation and constants depending on the protocol. Subsequent concrete analysis
shows that conjectured and provable parameter estimates can differ substantially, so real
systems should report the theorem or conjecture under which their claimed bit security is
computed~\cite{BlockTiwari2024ConcreteFRI}.

If one execution rejects a $\delta$-far word with probability at least $s(\delta)$, then $t$
independent query paths reduce the conditional acceptance probability to at most
\[
    \bigl(1-s(\delta)\bigr)^t.
\]
Concrete systems choose the code rate, folding arity, field, and number of queries together;
one should not estimate security from the number of queries alone.

\subsection{Compiling the oracle protocol with Merkle trees}

The IOP model gives the verifier ideal oracle access. A cryptographic proof realizes each
oracle using a Merkle commitment.

\subsubsection*{Commit phase with hashes}

For each evaluation vector
\[
    \mathbf f^{(i)}=\bigl(f^{(i)}(x)\bigr)_{x\in D_i},
\]
the prover builds a Merkle tree and publishes its root $R_i$. The transcript order is
\[
\begin{aligned}
R_0
&\longrightarrow \alpha_0=H(\mathsf{fri\mbox{-}challenge},R_0)\longrightarrow R_1\\
&\longrightarrow \alpha_1=H(\mathsf{fri\mbox{-}challenge},R_0,R_1)\longrightarrow R_2
\longrightarrow\cdots.
\end{aligned}
\]
The hash outputs are mapped into field elements. Domain-separation labels and the complete
statement must be included in the real transcript.

Collision resistance gives computational binding: after publishing $R_i$, the prover cannot
open one position to two different field values without producing a hash collision. This is
the hash-based analogue of the binding role played by a PCS commitment.

\subsubsection*{Grinding and Fiat--Shamir security accounting}

Replacing each verifier message by a hash does more than remove the network interaction: it
lets the prover evaluate the random oracle repeatedly before publishing a proof. If one complete
interactive transcript accepts falsely with probability $\varepsilon$, then an adversary trying
$Q$ independently varied transcript prefixes has the basic grinding opportunity suggested by
$Q\varepsilon$. For a $\kappa$-bit random oracle, a representative classical-ROM accounting for
FRI has the form~\cite{BlockGarretaKatzThalerTiwariZajac2023FRIFS}
\[
    \varepsilon_{\mathrm{FS}}
    \le
    Q\,\varepsilon_{\mathrm{rbr}}
    +O\!\left(\frac{Q^2}{2^\kappa}\right).
\]
Thus ``$\lambda$ bits interactively'' does not mean that an implementation may ignore the
attacker's hashing budget. For example, $2^b$ offline trials can consume roughly $b$ bits from
a naive $2^{-\lambda}$ target.

For an arbitrary many-round public-coin protocol, standard soundness alone need not prevent a
stronger attack that grinds one round at a time, retaining each favorable prefix. This is why the
lecture stresses \emph{round-by-round soundness}: from any doomed transcript prefix, the next
random challenge should leave the state doomed except with small probability. The lecture slides
predate the cited proof of this property for FRI; the correct modern conclusion is therefore not
that Fiat--Shamir FRI lacks a proof, but that non-interactive security relies on the dedicated
round-by-round/BCS analysis and its concrete parameters, rather than on the slogan ``replace every
challenge by a hash.''

\subsubsection*{Query phase with authentication paths}

After the final root or terminal polynomial has been sent, Fiat--Shamir derives the query
indices from the full transcript. For every queried value, the prover supplies a Merkle
authentication path. The verifier:
\begin{enumerate}[leftmargin=2em]
    \item reconstructs each root and checks that every opened value belongs to the previously
    committed oracle;
    \item checks the FRI folding equations along each path;
    \item checks the terminal polynomial and its degree bound.
\end{enumerate}

The IOP query complexity of binary FRI is $O(t\log d_0)$ field elements. With separate
binary Merkle paths, the authentication data along one path is naively $O(\log^2 N)$ hash
values, because the proof opens leaves in a sequence of progressively smaller trees. Batched
openings, Merkle multiproofs, larger folding arity, and larger tree arity reduce the concrete
overhead substantially.

\subsection{FRI is a low-degree test, not by itself a PCS}

It is common to say ``FRI commitment,'' but three layers should be distinguished:
\[
\begin{array}{c|c}
\text{Layer} & \text{Responsibility}\\
\hline
\text{Merkle tree} & \text{bind to an evaluation table}\\
\text{FRI IOPP} & \text{prove proximity of the table to a low-degree RS codeword}\\
\text{evaluation reduction} & \text{connect a claimed point value to that codeword}
\end{array}
\]

For example, to prove $p(z)=v$, define
\[
    q(X)=\frac{p(X)-v}{X-z}.
\]
If the claim is correct and $\deg p<d$, then $q$ is a polynomial of degree less than $d-1$.
On an evaluation domain $D$ not containing $z$, the verifier can check at sampled points that
\[
    p(x)-v=(x-z)q(x),
\]
while FRI proves that the committed $q$-table is close to a polynomial of the required degree.
In the simple unique-decoding interpretation, the distance must remain below roughly half the
relative minimum distance, about $(1-\rho)/2$, so that there is a unique nearby polynomial to
which the evaluation claim can refer. In that regime the elementary miss probability per query
is larger than $\rho$ and may contribute less than one bit per path, even when the raw
$\rho^t$ attack suggests $\log_2(1/\rho)$ bits. This is one reason FRI is often best understood
as providing \emph{list/proximity binding}: the Merkle root and FRI transcript constrain the
prover to a small set of nearby low-degree candidates, and the surrounding evaluation reduction
must select the candidate consistent with $p(z)=v$.
More refined protocols sample outside the evaluation domain to bind the prover more tightly
to a candidate polynomial. DEEP-FRI formalizes this ``domain extending for eliminating
pretenders'' idea and improves the soundness range while retaining linear prover arithmetic
and logarithmic verifier arithmetic~\cite{BenSassonGoldbergKoppartySaraf2020DEEPFRI}.

\subsection{Where FRI sits inside a STARK}

FRI does not prove arbitrary computation directly. In a STARK-style system the layers are:
\[
\boxed{
\begin{array}{c}
\text{execution trace}\longrightarrow\text{AIR constraints}\\
\longrightarrow\text{composition/quotient polynomial}\\
\longrightarrow\text{low-degree extension on a large domain}\\
\longrightarrow\text{FRI proximity proof}\\
\longrightarrow\text{Merkle commitments + Fiat--Shamir}.
\end{array}}
\]
AIR and the algebraic linking layer establish that the committed polynomial encodes a valid
execution if it is low degree. FRI supplies that final low-degree guarantee. Merkle trees make
the polynomial oracles succinctly accessible, and Fiat--Shamir removes interaction. This is
the hash-based route to scalable transparent arguments described by the STARK
construction~\cite{BenSassonBentovHoreshRiabzev2018STARK}.

\begin{remark}[Transparency and post-quantum security]
FRI requires no hidden algebraic trapdoor. Once compiled with Merkle trees, its cryptographic
assumption is principally the security of the hash function, which is why FRI-based systems
are described as transparent and plausibly post-quantum. Concrete post-quantum parameters
must still account for quantum attacks on the hash function and for the security analysis of
Fiat--Shamir in the quantum random-oracle model.
\end{remark}

\begin{remark}[Zero knowledge is an additional layer]
Plain FRI is not automatically zero knowledge. Its query answers reveal evaluations of the
committed codeword, and a Merkle root is binding but not an information-theoretic hiding
commitment. A zero-knowledge STARK randomizes or masks the underlying polynomial oracles
and proves that the masking preserves the required identities. FRI then proves low degree of
the masked objects. Thus transparency, proximity soundness, and zero knowledge are separate
properties that must each be established.
\end{remark}

\subsection{Complexity and design tradeoffs}

For the binary pedagogical protocol with $t$ query paths and $r=\log_2d_0$ folds:
\[
\begin{array}{c|c}
\text{Quantity} & \text{Asymptotic cost}\\
\hline
\text{folding arithmetic for the prover} & O(N)\\
\text{total committed field elements} & <2N\\
\text{verifier field queries} & O(t\log d_0)\\
\text{verifier field arithmetic} & O(t\log d_0)\\
\text{cryptographic setup} & \text{transparent; public hash function}\\
\text{proof size after Merkle compilation} & \text{polylogarithmic, but larger than KZG proofs}
\end{array}
\]

Higher-arity FRI decomposes
\[
    p(X)=\sum_{j=0}^{b-1}X^j p_j(X^b)
\]
and folds the $b$ coefficient polynomials using powers of a random challenge. It reduces the
number of rounds by a factor of $\log_2 b$ but opens more sibling values per round. This is a
typical proof-system tradeoff: larger arity shortens the recursion while increasing the local
work and the width of each query.

The enduring reason FRI is useful is not that any one cost is minimal. KZG has much smaller
openings and faster verification. FRI instead combines:
\begin{itemize}[leftmargin=2em]
    \item linear-time algebraic folding,
    \item logarithmic verifier arithmetic and query complexity,
    \item no trusted setup,
    \item hash-based commitments, and
    \item compatibility with large FFT-style execution traces.
\end{itemize}
Appendix~\ref{app:fri-example} works through every fold and one complete authenticated query
path over a small field.
